% ===================================================================
%  LENTELĖ Nr. 9 — Kalnas. Miškas.
%  Traduction 2026 d'apres texto/io, controlee sur texto/fr, texto/ru
%  et texto/cs. Consigne : voir texto/lt/10-lentele-01.tex.
%
%  LE TCHEQUE, TREIZIEME COLONNE, N'AVAIT PAS DE MER, et son tableau
%  10 avait montre ce qu'une langue fait alors : elle prete a la mer
%  les mots de la montagne — « útes » pour la falaise et pour l'ecueil,
%  « úskalí » tire de « skála ». LE LITUANIEN N'A PAS DE MONTAGNE, et
%  le tableau 9 montre la meme operation prise par l'autre bout.
%
%  IL N'A QU'UN MOT, « kalnas », POUR TOUT CE QUI MONTE. La butte de
%  vingt metres au bout d'un champ et le Mont-Blanc portent le meme
%  nom ; le point le plus haut du pays culmine a deux cent
%  quatre-vingt-quatorze metres. Tout le reste du tableau se fabrique
%  donc PAR COMPOSITION sur ce mot unique :
%
%    « kalnynas »     la chaine, le collectif de kalnas ;
%    « kalnagūbris »  la crete, l'echine du mont ;
%    « ugnikalnis »   le volcan, mot a mot le mont-de-feu ;
%    « kalnietis »    le montagnard.
%
%  ET « ugnikalnis » EST UN CALQUE, c'est-a-dire la troisieme voie que
%  le tcheque avait enseignee : ni garder le mot etranger, ni le
%  refaire de rien, mais le traduire piece par piece. Le lituanien la
%  prend ici pour une chose qu'il n'a jamais vue, et il garde en
%  revanche « lava » et « krateris » tels quels : la matiere
%  s'emprunte, la forme se calque.
%
%  LA SECONDE MOITIE DU TABLEAU RENVERSE LA PREMIERE. Le pays est
%  couvert de forets, et le vocabulaire y devient d'un coup dense et
%  entierement propre : « ąžuolas » le chene, « beržas » le bouleau,
%  « eglė » l'epicea, « pušis » le pin, « liepa » le tilleul, « guoba »
%  l'orme, « samanos » la mousse, « paparčiai » les fougeres,
%  « viržiai » la bruyere, « voverė » l'ecureuil, « genys » le pic,
%  « šarka » la pie, « elnias » le cerf, « stirna » le chevreuil,
%  « šernas » le sanglier, « kiškis » le lievre, « angis » la vipere,
%  « žaltys » la couleuvre, « skruzdėlynas » la fourmiliere.
%  UNE LANGUE EST RICHE LA OU SON PAYS L'EST.
%
%  ET LE CALENDRIER LITUANIEN EST FAIT DES MOTS DE CE LIVRET. Les
%  mois ne portent pas de noms latins : « liepa » juillet est le
%  tilleul, « birželis » juin le bouleau, « gegužė » mai le coucou,
%  « balandis » avril le pigeon, « rugpjūtis » aout la coupe du
%  seigle, « rugsėjis » septembre les semailles du seigle, « spalis »
%  octobre la chenevotte du lin, « lapkritis » novembre la chute des
%  feuilles. Les tableaux 3, 7 et 8 avaient nomme ces arbres, ces
%  oiseaux et ces travaux un a un ; le meme lexique fait ici l'annee
%  entiere. C'est pourquoi gravuri/korekti.json met en gras
%  « liepos » et non un « julius » : le mois EST le tilleul.
% ===================================================================

%%K t09-apar-1 apar
\VUsaut{4.0mm}
\VUtitre{15.0pt}{60}{LENTELĖ Nr. 9}
\VUsaut{3.0mm}

%%K t09-tit-1 sub
\VUcentre{12.0pt}{0}{\textit{Pirmoji scena.}}
\VUsaut{3.0mm}

%%K t09-tit-2 sub
\VUpk{11.4pt}{40}{Kalnas. --- Kurortas.}
\VUsaut{3.0mm}

%%K t09-c1-01-1 p
1. --- Jau aštuonios dienos, kai esu Prace. Negaliu išreikšti visos nuostabos, kurią
pajutau, kai atsidūriau prieš nuostabų Pirėnų \VUgras{kalnyną} \textsuperscript{(1)},
kurio \VUgras{ketera} \textsuperscript{(2)} mėlyname danguje atrodo kaip milžiniška
girlianda.

%%K t09-c1-02-1 p
2. --- Neįsivaizdavau, kad Pirėnų \VUgras{viršūnės} \textsuperscript{(3)} siekia tokį
didelį \VUgras{aukštį,} ir likau apstulbęs prieš puikius \VUgras{ledynus}
\textsuperscript{(5)} ir apsnigtas \VUgras{kalnų smailes} \textsuperscript{(4)}, nuo
kurių rieda tokios baisios \VUgras{lavinos.}

%%K t09-tit-3 sub
\VUsaut{3.0mm}
\VUpk{10.2pt}{40}{Kalnų peizažas.}
\VUsaut{3.0mm}

%%K t09-c1-03-1 p
3. --- Jau kitą dieną po atvykimo nuėjau prie senojo užgesusio \VUgras{ugnikalnio}
\textsuperscript{(6)}, kuris yra prie Praco. Ant nuogų ir plikų \VUgras{kraterio}
kraštų dar gerai matyti pėdsakai, kuriuos paliko \VUgras{lava,} tekėjusi prieš
kelis šimtmečius \VUgras{kalno šlaitu} \textsuperscript{(7)}. Pavyko man rasti ant
vieno iš \VUgras{šlaitų} keletą tos lavos nuolaužų.

%%K t09-c1-04-1 p
4. --- Pracas stovi pasienyje, ir \VUgras{tvirtovė} \textsuperscript{(8)}, kuri gina
\VUgras{tarpeklį} \textsuperscript{(27)}, skiriantį Prancūziją nuo Ispanijos, visai
primena tas senas Reino kranto pilis, kurios tarsi tyko \VUgras{grobio} nuo savo
uolos viršaus.

%%K t09-c1-05-1 p
5. --- Milžiniškas \VUgras{erelis} \textsuperscript{(9)}, kurio lizdas netoli
\VUgras{forto} \textsuperscript{(8)}, dažnai patraukia mano dėmesį.

Tas \VUgras{plėšrusis paukštis,} kurio \VUgras{snapas} stiprus, dažnai atneša savo
jaunikliams ėriuką ar ožiuką, kurį laiko savo \VUgras{nagais.} Tokie dideli jo
\VUgras{sparnai,} kad jis gali ilgai sklęsti su savo sunkia našta.

%%K t09-tit-4 sub
\VUsaut{3.0mm}
\VUpk{10.2pt}{40}{Ekskursantas.}
\VUsaut{3.0mm}

%%K t09-c1-06-1 p
6. --- Ten maloniausias pasivaikščiojimas yra išvyka prie „Kau“ krioklio
\textsuperscript{(10)}. \VUgras{Vandenpuolis} sūkuriuodamas krinta ir paskui dingsta
kaip gražus \VUgras{upelis} \VUgras{eglių giraitėje} \textsuperscript{(11)}, esančioje
į vakarus nuo miesto. Užkopėme per tą giraitę iki \VUgras{meteorologijos stoties}
\textsuperscript{(12)}, ir nuo \VUgras{apžvalgos aikštelės} viršaus pamatėme visą
srities \VUgras{panoramą.} \VUgras{Peizažas} nuostabus, ir \VUgras{gamta} tarsi
negailėdama beria kraštui visus turtus, kuriais disponuoja.

%%K t09-c1-07-1 p
7. --- Nusileisti pasinaudojome \VUgras{funikulieriumi} \textsuperscript{(13)} ir
sustojome mažoje \VUgras{stotelėje} prie senos \VUgras{feodalinės pilies}
\VUgras{griuvėsių} \textsuperscript{(14)}, kurioje kadaise gyveno Praco ponai.

%%K t09-c1-08-1 p
8. --- Vargšė pilis! Ką ji turbūt galvoja, matydama apačioje koketišką
\VUgras{kurortą} \textsuperscript{(15)}, kuris dabar yra madinga \VUgras{terminė
stotis}?

%%K t09-c1-08-2 p
\VUgras{Klimatas} toks malonus, \VUgras{mineralinis vanduo} toks veiksmingas, kad
\VUgras{gydymas,} taikomas \VUgras{sanatorijoje} \textsuperscript{(17)}, yra tikras
malonumas.

%%K t09-tit-5 sub
\VUsaut{3.0mm}
\VUpk{10.2pt}{40}{Malonus terminis miestas.}
\VUsaut{3.0mm}

%%K t09-c1-09-1 p
9. --- Beje, padaryta viskas, kad viešnagė būtų maloni. Buvo pastatytas didelis
\VUgras{viadukas} \textsuperscript{(16)}, kad būtų galima nutiesti geležinkelį;
\VUgras{terminės įstaigos} \VUgras{parkas} \textsuperscript{(18)}, kuriame
\VUgras{koncertuojama,} įrengtas žavingai. Kelios grakščios „\VUgras{vilos}“
\textsuperscript{(19)} pastatytos aplink
kazino \textsuperscript{(21)}, o labai gerai prižiūrimas \VUgras{plentas}
\textsuperscript{(22)} labai palengvina susisiekimą.

%%K t09-c1-10-1 p
10. --- Paprastas \VUgras{pylimas} \textsuperscript{(23)} skiria \VUgras{kelią}
\textsuperscript{(20)} nuo tikrojo \VUgras{slėnio} \textsuperscript{(25)}, kurio gale
guli grakštus mažas \VUgras{ežeras} \textsuperscript{(26)}, maitinantis skaidrų
\VUgras{upeliuką} \textsuperscript{(24)}.

%%K t09-c1-11-1 p
11. --- Kitoje slėnio pusėje eksploatuojama \VUgras{geležies kasykla}
\textsuperscript{(28)}. Kelis kartus ėjau pažiūrėti \VUgras{kalnakasių,} leidžiančių
į \VUgras{šachtą,} ir net įėjau į \VUgras{galeriją,} kurioje jie leidžia savo dieną,
kasdami \VUgras{rūdą,} kurioje yra \VUgras{metalo.}

%%K t09-c1-12-1 p
12. --- Prie kasyklos buvo pastatyta didelė \VUgras{liejykla,} kad tuoj pat būtų
panaudoti turtai, kurie iš jos iškasami. \VUgras{Aukštakrosnė} \textsuperscript{(29)},
kūrenama \VUgras{akmens anglimi,} kurios čia gausu, leidžia greitai ir pigiai gauti
\VUgras{ketaus.}

%%K t09-c1-13-1 p
13. --- Netoli nuo kasyklos kasamas \VUgras{akmens karjeras} \textsuperscript{(30)}.
Nei \VUgras{granito,} nei \VUgras{porfyro,} nei net \VUgras{smiltainio} senasis
\VUgras{akmenskaldys} nekerta savo \VUgras{kirtikliu,} o paprastus \VUgras{tašytus
akmenis} namams statyti.

%%K t09-c1-14-1 p
14. --- Viena gražiausių to krašto puošmenų yra \VUgras{eglės} \textsuperscript{(31)}.
Visur \VUgras{griovos} \textsuperscript{(32)} gilumoje, ant \VUgras{kalnų srauto}
\textsuperscript{(33)}, \VUgras{bedugnės} \textsuperscript{(34)} ar prarajos kranto, jų
plonos spyglės žaliai marguoja.

%%K t09-tit-6 sub
\VUsaut{3.0mm}
\VUpk{10.2pt}{40}{Ėjimas pėsčiomis.}
\VUsaut{3.0mm}

%%K t09-c1-15-1 p
15. --- Naudojuosi savo atostogomis, kad \VUgras{ilgai} vaikščiočiau.
\VUgras{Keliai} \textsuperscript{(35)}, kurie vingiuoja kalnu, leidžia nueiti,
nepaisant nuotolio, kelis \VUgras{kilometrus} ir dažnai kelis \VUgras{keturgubus
kilometrus.}

%%K t09-c1-15-2 p
\VUgras{Riboženkliai} \textsuperscript{(36)} ir \VUgras{rodyklių stulpai}
\textsuperscript{(37)} ženklina kelius, kuriuose dažnai sutinki jauną
\VUgras{kalnietį} \textsuperscript{(38)} su \VUgras{perpetine}
\textsuperscript{(40)} prie šono, nešantį \VUgras{švilpiką} \textsuperscript{(39)},
arba kokį \VUgras{čigoną} \textsuperscript{(41)}, kurio \VUgras{kailinukas}
\textsuperscript{(42)} keistai kontrastuoja su senu suplyšusiu \VUgras{gobtuvu}
\textsuperscript{(43)}. Jis veda grandine dresuotą \VUgras{lokį}
\textsuperscript{(44)}.

%%K t09-tit-7 sub
\VUpk{10.2pt}{40}{Kopimas į kalnus.}
\VUsaut{3.0mm}

%%K t09-c1-16-1 p
16. --- Beveik kasdien einu su \VUgras{vadovu} \textsuperscript{(48)} praktikuoti
\VUgras{kopimo,} ir esu labai išdidus, kai su \VUgras{kuprine}
\textsuperscript{(47)} ant nugaros ir „\VUgras{alpenštoku}“ rankoje, kaip tikras
\VUgras{turistas} \textsuperscript{(46)}, puolu \VUgras{uolas} \textsuperscript{(45)}.

%%K t09-c1-17-1 p
17. --- Nuo kalno viršaus lygumos gyventojai atrodo ne stambesni už skruzdes;
\VUgras{avių banda} \textsuperscript{(50)}, besigananti \VUgras{ganykloje}
\textsuperscript{(49)}, atrodo kaip vaikų žaislas. \VUgras{Pušis}
\textsuperscript{(51)} ar net \VUgras{medinis namelis} \textsuperscript{(52)} atrodo
vos kaip dėmelė žalumoje.

%%K t09-c1-18-1 p
18. --- Per tas išvykas dažnai sutinkame \VUgras{takelyje} \textsuperscript{(53)}
\VUgras{gemzių medžiotoją} \textsuperscript{(54)}, ant peties nešantį žvėrį, kurį ką
tik nudobė.

%%K t09-c1-19-1 p
19. --- Į savo herbariumą įdėjau labai įspūdingą \VUgras{paparčio}
\textsuperscript{(55)} šaką ir gražią rausvą \VUgras{viržio} \textsuperscript{(57)}
šakelę, kurias nuskyniau prie mažo \VUgras{urvo} \textsuperscript{(56)} angos per
paskutinį pasivaikščiojimą.

%%K t09-tit-8 sub
\VUsaut{3.0mm}
\VUcentre{12.0pt}{0}{\textit{Antroji scena.}}
\VUsaut{3.0mm}

%%K t09-tit-9 sub
\VUpk{11.4pt}{40}{Miškas. --- Medžioklė.}
\VUsaut{3.0mm}

%%K t09-c2-01-1 p
1. --- Vakar praleidau nuostabią dieną miške. Labai mane sudomino stropumas, kuris
viešpatauja tarp gyvūnų ir \VUgras{vabzdžių} \textsuperscript{(5)}, kurie jame gyvena.

%%K t09-c2-01-2 p
\VUgras{Voras} \textsuperscript{(1)}, kuris audžia savo \VUgras{tinklą}
\textsuperscript{(2)} tarp dviejų šakų, kad pagautų praskrendančią \VUgras{musę}
\textsuperscript{(3)}, \VUgras{sraigė} \textsuperscript{(4)}, vargais negalais nešanti
savo kiautą, elniavabalis, ieškantis savo grobio, stropi skruzdėlė, labai uoli prie
\VUgras{skruzdėlyno} \textsuperscript{(6)}, --- visi tie mažieji ėjo, atėjo, dirbo
nepaprastai stropiai.

%%K t09-c2-02-1 p
2. --- \VUgras{Gyvatė} \textsuperscript{(7)}, kurią iš pirmo žvilgsnio palaikiau
\VUgras{angimi,} bet kuri buvo ne kas kita kaip paprastas \VUgras{žaltys,} buvo
pritūpusi po medžio kamienu ir kartkartėmis iškišdavo savo ploną galvutę, kuri visada
atrodė pasirengusi kirsti. Priešais mane storas \VUgras{šliužas}
\textsuperscript{(8)} sunkiai šliaužė, prarijęs vieną iš skanių \VUgras{žemuogių}
\textsuperscript{(11)}, kurių ten gausiai auga.

%%K t09-c2-03-1 p
3. --- Žemė buvo išklota \VUgras{miško gėlėmis}: \VUgras{raktažolės}
\textsuperscript{(9)}, \VUgras{žiemės} \textsuperscript{(10)} ir daug kitų augalų,
kurių nepažinau, žydėjo tarp \VUgras{žolės kupstų} \textsuperscript{(12)}.

%%K t09-c2-04-1 p
4. --- Tame krašte \VUgras{trumas} beveik toks pat dažnas kaip \VUgras{grybas}
\textsuperscript{(14)}, ir eiguliui priklausantis \VUgras{paršelis}
\textsuperscript{(13)} gardžiuodamasis ieškojo, ar pavyks jam rasti kelių.

%%K t09-tit-10 sub
\VUsaut{3.0mm}
\VUpk{10.2pt}{40}{Brakonieriavimas.}
\VUsaut{3.0mm}

%%K t09-c2-05-1 p
5. --- Mačiau vienos scenos \VUgras{pabaigą,} kuri mane labai palinksmino.
Padedamas savo \VUgras{trumpakojo šuns} \textsuperscript{(15)} uoslės,
\VUgras{eigulys} \textsuperscript{(16)} ką tik užklupo \VUgras{brakonierių}
\textsuperscript{(22)} tą akimirką, kai šis rengėsi išsinešti \VUgras{žvėrieną}
\textsuperscript{(17)}, kurią buvo nudobęs. Verčiau palikdamas savo grobį, negu
leisdamasis suimamas, brakonierius pabėgo, o \VUgras{kiškis}
\textsuperscript{(18)}, \VUgras{elnė} \textsuperscript{(19)}, \VUgras{stirna}
\textsuperscript{(20)} ir \VUgras{šernas} \textsuperscript{(21)} liko gulėti ant
žolės, džiugindami eigulį.

%%K t09-c2-06-1 p
6. --- Nuostabi medžių, kuriuos apima miškas, įvairovė. \VUgras{Bukai}
\textsuperscript{(23)} ir \VUgras{beržai} \textsuperscript{(26)}, \VUgras{pušys,}
kurios duoda \VUgras{sakų} \textsuperscript{(27)}, \VUgras{ąžuolai}
\textsuperscript{(29)} ir dar daug kitų maišo savo šakas.

%%K t09-c2-06-2 p
\VUgras{Gebenė} \textsuperscript{(24)}, kuri glaudžiasi prie aukštų medžių kamienų,
nudažo \VUgras{girią} \textsuperscript{(25)} spindinčia žaluma, kuri maloniai
susilieja su šviesesniu ir blyškiai žaliu \VUgras{samanų} \textsuperscript{(31)}
atspalviu, kuriose \VUgras{žalioji meleta} \textsuperscript{(30)} ieško vabzdžių.

%%K t09-tit-11 sub
\VUsaut{3.0mm}
\VUpk{10.2pt}{40}{Miško peizažas.}
\VUsaut{3.0mm}

%%K t09-c2-07-1 p
7. --- Senieji ąžuolai mane sužavėjo. Jų storos susisukusios \VUgras{šaknys}
\textsuperscript{(32)}, pusiau išlindusios iš žemės, suteikia jiems puikią stiprybės
ir galios išvaizdą, ir klausiu savęs, kaip \VUgras{savininkas}
\textsuperscript{(35)} gali ryžtis juos nukirsdinti ir atiduoti \VUgras{pjūklui}
\textsuperscript{(34)}, kuriuo dirba \VUgras{medkirtis} \textsuperscript{(33)}. Vargšai
\VUgras{kamienai} \textsuperscript{(36)}, nuplėštos \VUgras{žievės}
\textsuperscript{(37)} arba skaldomi \VUgras{kirviu} \textsuperscript{(38)}, kurį kelia
\VUgras{dienininkas} \textsuperscript{(39)}, mane liūdina.

%%K t09-c2-08-1 p
8. --- Netoliese senas \VUgras{anglių degėjas} \textsuperscript{(41)} savo
\VUgras{kastuvu} sukrovė anglių \VUgras{krūvą} \textsuperscript{(42)}, o senutė
nešėsi \VUgras{glėbį} \textsuperscript{(40)} \VUgras{sausų malkų} iš nukirstų medžių
šakelių.

%%K t09-tit-12 sub
\VUsaut{3.0mm}
\VUpk{10.2pt}{40}{Medžioklė.}
\VUsaut{3.0mm}

%%K t09-c2-09-1 p
9. --- Norėjau eiti kelioms akimirkoms pailsėti į \VUgras{eigulio namą}
\textsuperscript{(46)}, bet, kai priėjau \VUgras{ežerėlį} \textsuperscript{(43)},
kuriame plaukioja graži \VUgras{gulbė} \textsuperscript{(45)}, pastebėjau, kad
neseniai buvęs \VUgras{potvynis} \textsuperscript{(44)} pavertė kelią tikra
\VUgras{pelke.} Turėjau apsigręžti ir, eidamas pro \VUgras{kedrą}
\textsuperscript{(49)}, \VUgras{tuopas} \textsuperscript{(48)} ir \VUgras{guobą}
\textsuperscript{(47)}, kurios stovi miško pakraštyje, pamačiau iš \VUgras{jaunuolyno}
\textsuperscript{(54)} išnyrant puikų \VUgras{elnią} \textsuperscript{(50)},
persekiojamą atkaklios \VUgras{medžioklinių šunų rujos} \textsuperscript{(51)},
kurią kurstė ir \VUgras{ragas} \textsuperscript{(53)}, pučiamas \VUgras{raitų
medžiotojų} \textsuperscript{(52)}. Vargšas žvėris buvo visai išsekęs. Jis atėjo
kelių žingsnių
atstumu nuo manęs mirdamas parkristi.

%%K t09-c2-10-1 p
10. --- Eigulio sūnus, kurį buvo pritraukęs \VUgras{varymo medžioklės} triukšmas,
išėjo iš namų, ir mudu grįžome į mišką. Jis buvo matęs \VUgras{šarką}
\textsuperscript{(56)} ant seno ąžuolo šakos. Jis manė, kad jos lizdas turbūt
paslėptas \VUgras{lapijoje} \textsuperscript{(58)}, ir norėjo lizdą pačiupti.

%%K t09-c2-10-2 p
Vikrus kaip \VUgras{voverė} \textsuperscript{(61)}, jis lipo iš \VUgras{šakos}
\textsuperscript{(55)} į šaką, bet nieko nerado ir tepavyko jam pakelti skristi seną
\VUgras{varną} \textsuperscript{(60)}, tupėjusią ant \VUgras{šakos}
\textsuperscript{(57)}.

\VUsaut{4.0mm}
\VUfilet{22mm}
